\documentclass[12pt]{article}
\usepackage{enumitem}
\usepackage{amsmath}
\usepackage{geometry}
\geometry{margin=1in}

% Reduce spacing after \maketitle
\usepackage{etoolbox}
\makeatletter
\renewcommand{\maketitle}{
  \begin{center}
    {\Large \@title \par}
    \vskip 0.5em % space between title and date (reduce this as needed)
    {\@date}
  \end{center}
  \vskip 0em % space after date before text
}
\makeatother

\title{\Large List of Possible Future Projects}
\date{\today}

% Define shortcuts for bullet points
\newcommand{\point}[1]{\vspace{1.2em}\noindent$\bullet$ \textbf{#1}\par}
\newcommand{\subpoint}[1]{\vspace{0.3em}\noindent\hspace{1em}-- #1\par}

\begin{document}

\maketitle

\point{Considering nucleation of SQM seed}
\subpoint{Suppressed exponentially. Energy required to nucleate $\propto s R^3 + \sigma R^2$. The volume term is favored but the surface term is not, and when $R$ is small, the bubble isn't favored. So there's a critical radius $R_c$ above which bubbles can expand. The rate for a finite number of $s$ quarks to be created at the same time near each other is supressed exponentially by the number of strange quarks. }

\point{Use more sophisticated models in different directions}
\subpoint{Consider RMF models for NM (already done)}
\subpoint{Use RMF-like model for QM.}

\point{Generic picture \& post merger signals}
\subpoint{NM $\rightarrow$ SQM transition will most likely be observed in post-merger signals once the SQM seed is nucleated in the high temperature environment of the inspiral event due to the fact of small $v_{N\to Q}\sim 10^1$(m/s)}

\point{V-QCD and gauge/gravity modeling}
\subpoint{Modeling in the non-perturbative region of QCD.}

\point{Isospin and Urca processes}
\subpoint{Consider isospin and Urca processes as well. Leading to two coupled differential equations}
\subpoint{For two-flavor QM, Urca process is the only weak process. }

\point{Processes ahead of the phase boundary}
\subpoint{Are there any advanced processes happening ahead of the phase boundary?}

\point{Breaking color superconductivity}
\subpoint{Breaking color superconductivity Cooper pairs.}

\point{Equation of State as a function of $T$}
\subpoint{Consider EoS with non-zero temperature.}

\point{SQM to NM transition with hyperons}
\subpoint{Include hyperons in NM for the SQM to NM transition.}

\point{Pressure assumption validity}
\subpoint{$P_{Q^*} = P_Q = P_b$ is not a trivial assumption; it depends on the scale of the transition region. If it's larger, then maybe pressure can't be treated as constant.}

\point{Assumptions on the direction of Kaon-ness isospin diffusion}
\subpoint{Is it good to assume $a_N > a_{Q^*}$? When does this assumption break down?}

\point{Improvement on Diffusion coefficient.}
\subpoint{The $D_Q$ is a crude / sociologically acceptable approximation, indicate that calculating it using something like $\chi$EFT is worthwhile}

\newpage

\subsection*{Outline: NM $\to$ SQM Conversion in Neutron-Star Mergers}

\paragraph{Goal.}
\begin{itemize}
    \item Construct a simplified, microphysics-based theory of NM $\to$ SQM conversion.
    \item Derive conservative upper bounds on conversion speed.
    \item Put in NSM conditions (two phases: inspiral/merger and post-merger)
    \item Simplified enough/capture the dominant contribution.
\end{itemize}

\paragraph{Two conversion regimes.}
We distinguish two physically distinct regimes for NM $\to$ SQM conversion:
\begin{itemize}
    \item \textbf{Fast (hydrodynamic) mode}
    \begin{itemize}
        \item Operates when hydrodynamic driving is strong.
        \item Relevant during inspiral and merger phases.
        \item Dominated by compression, shocks, and turbulence.
        \item Coll's condition is satisfied.
        \item Diffusion-limited models are not applicable.
    \end{itemize}
    
    \item \textbf{Slow (diffusion-limited) mode}
    \begin{itemize}
        \item Operates when hydrodynamic driving becomes ineffective.
        \item Relevant in the post-merger, quasi-hydrostatic remnant.
        \item Coll's condition fails.
        \item Conversion proceeds via diffusion and weak interactions.
    \end{itemize}
\end{itemize}

\paragraph{Fast mode and nucleation.}
\begin{itemize}
    \item The fast mode may enable the nucleation of strange-quark-matter seeds.
    \item Such nucleation requires extreme thermodynamic conditions.
    \item Modeling nucleation in this regime might lie beyond the scope of the present work.
    \item A phenomenological treatment of strange-nugget nucleation is left for future study.
\end{itemize}

\paragraph{Slow mode and front propagation.}
\begin{itemize}
    \item Once hydrodynamic driving ceases, conversion must proceed in the slow mode.
    \item In this regime, the propagation speed of the conversion front is controlled by:
    \begin{itemize}
        \item diffusion,
        \item weak interaction rates.
    \end{itemize}
    \item The present work focuses on this regime.
\end{itemize}

\paragraph{Leptonic \& neutrino \& intermediate udQM matter effects.}
\begin{itemize}
    \item Leptonic (Urca) processes, neutrino transport, and intermediate udQM matter may modify weak equilibration.
    \item Their impact on the diffusion-limited front speed should be further assessed.
    \item This determines whether such effects are essential or subleading under NSM conditions.
\end{itemize}



\end{document}